\documentclass[10pt]{article}

% --- Packages ---
\usepackage[a4paper,margin=0.6in,top=0.4in,bottom=0.3in,left=0.6in,right=0.6in]{geometry}
\usepackage{enumitem}
\usepackage{hyperref}
\usepackage{titlesec}

% --- Formatting ---
\pagenumbering{gobble}
% \titleformat{\section}{\bfseries}{}{0em}{}[\titlerule]
\titlespacing{\section}{0pt}{1ex plus 0.4ex minus 0.2ex}{0.9ex plus 0.2ex minus 0.2ex}
\setlist[itemize]{leftmargin=1em, rightmargin = 1em,itemsep=0.7pt, topsep=1pt}
\titleformat{\section}{\bfseries}{}{0em}{}[\titlerule]
\renewcommand{\baselinestretch}{0.92}
% --- Document ---
\begin{document}

\begin{center}
    {\Large \textbf{Joseph M. Ryan}} \\[1pt]
    \href{mailto:joseph.ryan@duke.edu}{joseph.ryan@duke.edu}    Languages: English and French (native), Spanish (working)\\
    Background: Winchester College alumnus
\end{center}


% --- Education ---
\section*{Education}
\noindent\textbf{Duke University}, Ph.D. Candidate in Physics (Quantum Computing) at the Duke Quantum Center. \hfill 2022--present \\
\indent  Advisor: Crystal Noel,  Theory Collaborators: Stephen Teitsworth. \\
\noindent
\textbf{University of Michigan}, B.S. in Honors Physics \& Mathematics, High Honors. \hfill 2022 \\
\indent Recipient of William L. Williams Award for thesis on QCD color entanglement (LHCb).

% --- Research Summary ---
\section*{Research Summary}
\noindent Physicist specializing in \textbf{quantum stochastic processes.} 

% --- Research Experience ---
\section*{Research Experience}
\noindent\textbf{Duke Quantum Center (Noel Lab)} \hfill 2022--present
\begin{itemize}
    \item Built and operated a trapped-ion quantum computer
\end{itemize}



\noindent\textbf{Univ. of Michigan -- Aidala Group (LHCb Experiment at CERN)} \hfill 2021--22
\begin{itemize}
    \item Designed statistical methods to detect TMD factorization breaking due to color entanglement in proton collisions; analyzed correlated stochastic parton dynamics.
\end{itemize}

% --- Publications ---
\section*{Publications}
\begin{itemize}
    \item J. M. Ryan \textit{et al.}, ``Experimental Measurement of Quantum First-Passage-Time Distributions.''  
    \href{https://journals.aps.org/prresearch/abstract/10.1103/p3pl-1m6n}{Physical Review Research 8, L022025, (2026)}. Featured in \textit{Quantum Zeitgeist}.
    \item J. M. Ryan \textit{et al.}, ``Quantum-classical crossover in noisy monitored oscillators.'' 
    \href{https://arxiv.org/abs/2608.08267}{arXiv:2608.08267v1 (2026)}.
\end{itemize}

% --- Talks ---
\section*{Selected Talks}
\begin{itemize}
    \item APS DAMOP (2025) -- \textit{Noise-Driven Quantum First-Passage-Time Distributions in Ion Traps}
    \item NSF Robust Quantum Computing, Princeton (2024) -- \textit{Quantum Brownian Motion and First-Passage Times}.
    \item APS March Meeting (2024) -- \textit{Noise-Driven Escape Times in Trapped Ion Systems}.
\end{itemize}


% --- Work Experience ---
\section*{Work Experience}
\noindent\textbf{OTC Fellow, Duke Office of Translation \& Commercialization} \hfill 2025-present
\begin{itemize}
    \item Supporting Duke startups through market analysis and go-to-market strategy development.
\end{itemize}

\noindent\textbf{Advanced Physics Laboratory, Univ. of Michigan} \hfill 2020
\begin{itemize}
    \item Teaching Assistant.  Designed and taught advanced atomic, optical and nuclear physics experiments.
\end{itemize}

% 
    
\noindent\textbf{Michigan Hyperloop Team} \hfill 2018--19
\begin{itemize}
    \item Business director. Raised and managed funds for a multidisciplinary engineering team for hyperloop.  
\end{itemize}

\noindent\textbf{Airbus} \hfill 2017
\begin{itemize}
    \item Metrology Intern. Developed precision measurement protocols to ensure quality and accuracy in wing production.  
\end{itemize}

% \noindent\textbf{Airbus (metrology department, Broughton UK)} \hfill Summer 2017
% \begin{itemize}
%     \item Intern.  Developed automated calibration routines to improve precision of factory robots.
% \end{itemize}




% --- Skills ---
\section*{Skills}
\begin{itemize}
    \item \textbf{Stochastic Modeling:} Quantum/classical stochastic processes, first-passage times
\end{itemize}

\end{document}
